% Experiments section -- AAAI-27 (two-column, US Letter).
%
% PREAMBLE: this file adds ONE package to whatever the AAAI kit loads:
%     \usepackage{booktabs}
% booktabs only defines \toprule/\midrule/\bottomrule/\cmidrule and touches nothing in the
% page layout, so it is compatible with the kit's restrictions. If your reviewers' build does
% not have it, replace \toprule \midrule \bottomrule with \hline and \cmidrule(lr){a-b}
% with \cline{a-b}; nothing else changes.
%
% AAAI CONFORMANCE, checked against the submission rules:
%   - no geometry, no hyperref, no fontenc -- none are used here
%   - no \vspace, \hspace or any macro that alters margins, font size or line spacing
%   - no \includegraphics, so no trim/clip and no image-format question
%   - no page numbers, headers or footers set here
%   - two-column assumed: tab:main and tab:released span both columns via table*;
%     tab:budget, tab:timestep and tab:scaling are sized to fit one 3.3in column
%   - \small inside a float is ordinary table sizing, not a layout change
%
% FLOAT PLACEMENT: table* uses [!t] and table uses [!tb]. The "!" tells LaTeX to
% ignore \topfraction/\totalnumber/\dbltopfraction for that float, which is what
% makes five floats queue up and get flushed together at the \FloatBarrier that
% aaai2027.sty:422 puts before the bibliography. "b" is not available to table* --
% bottom placement for a spanning float needs stfloats, which the kit forbids.
% These are placement specifiers, not style changes, so they are permitted.
%
% PAGE BUDGET: trimmed to two floats for AAAI-27's 7 content pages. Three tables were cut
% and their numbers put in prose -- tab:scaling and tab:budget as sentences, tab:released
% to the supplementary document (supplementary.tex). tab:main and tab:timestep are kept
% because they carry the claims. Nothing was dropped: every number still appears.
%
% CONFIGURATION: every number is the model exactly as Method describes it, at a single
% training seed. The paper does not describe or report any post-hoc refinement stage, so
% none of these figures include one -- do not paste in numbers from the internal records
% that do, because they are a different configuration and would not match Method.
%
% MINIMUM REGION SIZE (added after the first draft, see \paragraph{Minimum region size}):
% tab:main is now reported WITH the rule at f = 0.01, applied to every model including the
% baselines. tab:timestep and the scaling paragraph are WITHOUT it, and say so, because they
% are within-model comparisons. The without-rule numbers for tab:main, if they are ever
% needed: ours .8421/.7545/.6410 (CLEVR), .8580/.7885/.8173 (Tetrominoes),
% .6974/.8177/.7628 (Multi-dSprites); best-baseline CLEVR foreground IoU .4219.
%
% DO NOT put akornC_* numbers in this table. Those files are the superseded AKOrN (attn) run
% at the repository's default seed 1234 (CLEVR fg_ari .4688); the seed-0 rerun this paper
% reports is in the s0_* files (CLEVR .6027/.3327/.3750). Reading akornC_tetrominoes would
% put foreground IoU .7621 in the best-baseline row and understate the margin.

\section{Experiments}
\label{sec:experiments}

\subsection{Protocol}
\label{sec:protocol}

All models are evaluated by one function on one representation. Ground truth is quantised to a
$64\times64$ patch grid and scored there; a prediction is \emph{never} upsampled to pixels to be
compared. A prediction coarser than the grid is lifted by an exact \texttt{repeat}, and one finer
than the grid is reduced by patch majority vote, so every method is scored on the same object. The
ceiling is therefore $1.0$ by construction, and these numbers are not directly comparable with
pixel-protocol figures reported elsewhere.

We report three metrics and never combine them into a single score: foreground adjusted Rand index
(\textsc{fg-ari}), foreground IoU (\textsc{fg-iou}) and matched-object IoU (\textsc{obj-iou}).
Background is label $0$ on both sides, and the foreground-only metrics exclude patches the ground
truth calls background.

\paragraph{Matched exposure.}
Every model is trained on the \emph{same} $10{,}000$ CLEVR images for $10$ epochs, i.e.\
$100{,}032$ exposures. Where a baseline's published recipe differs, only its epoch count is
changed; architecture and optimiser are left at the authors' values. Tetrominoes and
Multi-dSprites are never trained on: for every row in Table~\ref{tab:main} those two columns are
zero-shot transfer from the CLEVR-trained model.

\paragraph{Readouts that require a cluster count.}
Several baselines produce a continuous embedding and obtain masks by clustering it. We found that
the reference implementations of CAE, CAE\texttt{++}, CtCAE and SynCx set the number of clusters
from the \emph{per-image ground-truth object count},
\[
  k_i \;=\; \bigl| \mathrm{unique}(y_i) \bigr| ,
\]
which makes their published figures oracle-assisted. For the comparison to be meaningful we
removed this: $k$ is a per-dataset constant selected on a validation split, never read from the
evaluated image. Cluster identities are arbitrary, so the mapping of one cluster to background is
likewise fixed by a rule chosen on validation. No hyperparameter, threshold or rule is selected on
the reported split.

\paragraph{Minimum region size.}
\label{sec:minsize}
Our readout emits a sizeable number of single-patch regions: $22.8\%$ of all predicted regions on
CLEVR validation, $27.0\%$ on Tetrominoes and $56.2\%$ on Multi-dSprites, where the \emph{median}
predicted region is one patch. They arise after the synchrony grouping, when the greedy merge
splits a spatially scattered synchronous component into pieces, and they inflate the object count
without materially affecting the partition. We therefore discard any predicted region smaller
than a fraction $f$ of that image's predicted foreground, reassigning it to background. The rule
is scale-free, reads no ground truth, and $f$ is a single constant fixed on CLEVR validation as
the largest value whose FG-ARI loss stays under $0.005$ and whose matched-object IoU loss stays
under $0.01$, giving $f=0.01$. The same $f$ is applied unchanged to Tetrominoes and
Multi-dSprites --- no target-domain data is used to set it --- and to every baseline.

\subsection{Generalisation across datasets}
\label{sec:generalisation}

Table~\ref{tab:main} is the central result. Every model sees only CLEVR during training; the
Tetrominoes and Multi-dSprites columns measure transfer.

\begin{table*}[!t]
\centering
\small
\setlength{\tabcolsep}{4pt}
\begin{tabular}{l ccc ccc ccc}
\toprule
& \multicolumn{3}{c}{CLEVR (in-domain)} & \multicolumn{3}{c}{Tetrominoes (zero-shot)} & \multicolumn{3}{c}{Multi-dSprites (zero-shot)} \\
\cmidrule(lr){2-4}\cmidrule(lr){5-7}\cmidrule(lr){8-10}
Model & \textsc{ari} & \textsc{f-iou} & \textsc{o-iou} & \textsc{ari} & \textsc{f-iou} & \textsc{o-iou} & \textsc{ari} & \textsc{f-iou} & \textsc{o-iou} \\
\midrule
\textbf{Ours} & \textbf{.8408} & \textbf{.7528} & \textbf{.6387} & \textbf{.8568} & \textbf{.7950} & \textbf{.8170} & \textbf{.6966} & \textbf{.8190} & \textbf{.7625} \\
\midrule
AKOrN (attn) $+$ up-tiling & .6326 & .3468 & .3636 & .5484 & .7230 & .4167 & .5505 & .6539 & .6348 \\
AKOrN (conv)               & .6046 & .4219 & .2769 & .7202 & .7540 & .5913 & .4055 & .6834 & .5434 \\
AKOrN (attn)               & .6027 & .3327 & .3750 & .5676 & .7383 & .4687 & .4343 & .6019 & .5751 \\
ItrSA                      & .5794 & .2787 & .3853 & .6319 & .7137 & .5687 & .4817 & .2489 & .6093 \\
CAE\texttt{++}             & .2838 & .2193 & .0882 & .2715 & .2452 & .1901 & .2648 & .1048 & .1125 \\
Slot Attention             & .2872 & .2220 & .0801 & .1739 & .2192 & .1071 & .1930 & .0876 & .0680 \\
CtCAE                      & .2357 & .2148 & .0964 & .2681 & .2420 & .2034 & .3079 & .1158 & .1609 \\
CAE                        & .1091 & .2008 & .0634 & .1526 & .2305 & .1192 & .1200 & .1038 & .0630 \\
\midrule
Best baseline, per cell    & .6326 & .4267 & .3853 & .7202 & .7540 & .5913 & .5505 & .6834 & .6348 \\
\textbf{Our margin}        & \textbf{+.2082} & \textbf{+.3261} & \textbf{+.2534} & \textbf{+.1366} & \textbf{+.0410} & \textbf{+.2257} & \textbf{+.1461} & \textbf{+.1356} & \textbf{+.1277} \\
\bottomrule
\end{tabular}
\caption{%
All models trained on the same $10{,}000$ CLEVR images for $10$ epochs and evaluated on $320$
held-out images per dataset by the same function. Tetrominoes and Multi-dSprites are zero-shot.
\textsc{ari}\,/\,\textsc{f-iou}\,/\,\textsc{o-iou} are foreground ARI, foreground IoU and
matched-object IoU on the $64\times64$ patch grid. Every model, ours included, has the
minimum-region-size rule of Section~\ref{sec:minsize} applied at the single value $f=0.01$ fixed
on CLEVR validation. It leaves the baseline rows unchanged to three decimals except AKOrN (conv)
on CLEVR foreground IoU, which improves from $.4219$ to $.4267$; the penultimate row uses the
better of the two for every baseline cell we measured. That row takes the best baseline \emph{in
each cell separately}, which no single baseline attains; our margin is measured against that
composite. Single training seed throughout.}
\label{tab:main}
\end{table*}

Our readout is ahead in all nine cells. The comparison is deliberately pessimistic: the
\emph{best baseline per cell} is a composite of three different architectures, and no single
competitor achieves that row. It is also notable that \emph{which} baseline comes second depends
on the dataset --- AKOrN (conv) on Tetrominoes, AKOrN (attn)$+$up-tiling on Multi-dSprites, ItrSA
on CLEVR matched-object IoU --- whereas our model is first everywhere.

\paragraph{Where the margin is thin.}
The smallest margin in the table is Tetrominoes foreground IoU, $+0.0410$. Every other margin lies
between $0.128$ and $0.326$. We flag this cell explicitly rather than quoting only the aggregate.

\paragraph{Predicted foreground.}
A single diagnostic separates grouping from tiling. Against a true foreground fraction of
$0.2115$ on CLEVR and $0.0943$ on Multi-dSprites, our readout predicts $0.1988$ and $0.0827$ ---
the only method that does not over-claim --- while baselines predict between $0.41$ and $0.86$ on
CLEVR. A model that labels most of the image as object can score well on a foreground-only
partition measure while failing the IoU metrics, which is what the \textsc{f-iou} column shows.

\subsection{Are the baselines handicapped?}
\label{sec:fairness}

Two objections to Table~\ref{tab:main} deserve measurements rather than argument: that the
baselines are evaluated away from their published setting, and that the matched budget starves
them.

\paragraph{Released weights.}
Only Slot Attention ships trained weights; the AKOrN, CAE and SynCx repositories release data but
no checkpoints. Evaluated directly, at roughly $320\times$ our exposure budget, its released model
scores $.3747/.2219/.1165$ on CLEVR, $.4742/.2381/.2355$ on Tetrominoes and $.3205/.0980/.1204$
on Multi-dSprites. We remain ahead in all nine cells, by $0.377$ to $0.720$, so the comparison
does not rest on our retraining of it. The extra data also helps it far more \emph{off} its
training distribution than on it --- worth $+.0875$ \textsc{fg-ari} on CLEVR but $+.3003$ on
Tetrominoes and $+.1275$ on Multi-dSprites --- so a data-efficiency conclusion drawn from an
in-domain column alone would not transfer.

\paragraph{Budget.}
Comparing each baseline against its own published number, by its own evaluation code, isolates
the effect of the shortened schedule. Two of the six are \emph{not} budget-limited: AKOrN (conv)
reaches $60.5$ at $1/210$ of its schedule against a published $59.0\pm4.3$, and CAE reaches $10.3$
at $1/64$ against a published $4.0$. For those two the margin in Table~\ref{tab:main} is not an
artefact of a short schedule. The other four are limited --- AKOrN (attn) $60.4$ against
$91.0\pm0.5$, ItrSA $58.1$ against $82.5\pm1.7$, CAE\texttt{++} $16.7$ against $27.0$, and CtCAE
$12.0$ against $54.0$ --- and against those the result should be read as a data-efficiency claim.
We state both readings rather than only the favourable one. Full figures are in the supplementary
document.

\subsection{Rollout length}
\label{sec:timestep}

The readout integrates spiking activity over $T$ simulation steps and reads the second half of the
window. Table~\ref{tab:timestep} varies $T$ with the model fixed.

\begin{table}[!tb]
\centering
\small
\begin{tabular}{r ccc r}
\toprule
$T$ & \textsc{fg-ari} & \textsc{fg-iou} & \textsc{obj-iou} & $\Delta$\,\textsc{fg-ari} \\
\midrule
$1024$ & .8421 & .7545 & .6410 & --- \\
$512$  & .8402 & \textbf{.7594} & .6340 & $-.0019$ \\
$256$  & .8159 & .7358 & .6084 & $-.0262$ \\
$128$  & .5796 & .4988 & .3985 & $-.2625$ \\
\bottomrule
\end{tabular}
\caption{CLEVR, model fixed, readout window $T/2$. Halving the rollout is free; a quarter costs
$0.02$--$0.03$; an eighth fails. This table and the scaling figures below are reported
\emph{without} the minimum-region-size rule, because both are within-model comparisons at a
fixed readout and the rule would shift every row by the same small amount.}
\label{tab:timestep}
\end{table}

Halving the rollout costs $0.0019$ \textsc{fg-ari} --- below any difference we can resolve at one
seed --- and \emph{improves} foreground IoU. Performance degrades gently to $T=256$ and collapses
at $T=128$, where the model also under-segments, predicting a foreground fraction of $0.1148$
against a truth of $0.2115$. We therefore use $T=512$ as the default operating point and $T=1024$
for the reported numbers. Shorter windows also fragment more: mean predicted regions per image
rise from $8.0$ at $T=1024$ to $13.3$ at $T=256$, against $6.36$ true objects.

\paragraph{Scaling with training data.}
Continuing training on additional unseen CLEVR images, in steps of $1{,}000$ from $10{,}000$ to
$18{,}000$, separates the three metrics. \textsc{fg-ari} runs
$.8421$, $.8426$, $.8449$, $.8460$, $.8469$, $.8459$, $.8456$, $.8470$, $.8461$: it gains
$.0028$ over the first two steps and then stays inside a band of $.0021$ for the remaining six,
with no trend. Foreground IoU runs $.7545$ to $.7755$ and matched-object IoU $.6410$ to $.6579$,
both still rising at $18{,}000$ and neither turning. More in-domain data therefore stops improving
\emph{how} patches are grouped at around $12{,}000$ images while it is still improving \emph{where
the boundaries fall}. Mean predicted regions hold at $7.7$--$8.0$ throughout, against $6.36$ true
objects, so the extra data sharpens the partition without changing how much it splits. The
\textsc{fg-ari} band is one seed wide and we read it as flat rather than converged; the two IoU
curves have not been run far enough to find their own ceiling.

\paragraph{In-domain data is not the most useful data.}
The same $5{,}000$ additional images buy more than twice as much when they come from a different
dataset. Continuing on CLEVR raises CLEVR \textsc{fg-ari} by $+.0038$; continuing on Tetrominoes
instead raises it by $+.0084$, a $2.2\times$ larger gain at an identical budget. Together with
the transfer result in Table~\ref{tab:main}, this argues that the model is not fitting a domain
but learning a grouping rule that the extra renders refine.

\subsection{Spike timing, measured}
\label{sec:timing}

The two measurements summarised in Section 3.2, in full. Both run on the first $40$ images of
the CLEVR reserved split, on the same rollouts, so they are directly comparable.

\paragraph{Validity of the ablation.}
The affinity $Q$ is recomputed here from \texttt{component\_spikes} rather than taken from the
model, so that a modified spike trace can be fed through the same arithmetic. On the unshifted
control the recomputed $Q$ agrees with the model's own to $6.7\times10^{-6}$ maximum absolute
difference over all $4096^2$ entries, so the ablations differ from the control in exactly one
thing.

\paragraph{Separability.}
Ground-truth labels are used only to sort patch pairs into three groups and never to set a
threshold or a count. Averaged over the $40$ images, the affinity is $0.891\pm0.111$ for pairs
inside one object, $0.047\pm0.048$ for pairs in different objects, and $0.099\pm0.106$ for
object--background pairs. The within-object value is nineteen times the across-object value, and
the across-object value sits near the floor of the measure.

\paragraph{Per-patch circular shift.}
Each patch's four component traces are rolled by one offset drawn uniformly over the readout
window, independently per patch. Rate, waveform and within-patch structure survive; only the
alignment between patches is destroyed. The graph, the foreground threshold, the component
floor, the background heuristic and the greedy merge are untouched.

\begin{table}[!tb]
\centering
\small
\begin{tabular}{l ccc r}
\toprule
& \textsc{fg-ari} & \textsc{fg-iou} & \textsc{obj-iou} & regions \\
\midrule
control            & .8648 & .8028 & .6925 & 7.8 \\
per-patch shift    & .0735 & .7358 & .0325 & 562.3 \\
gates, not spikes  & .8648 & .8028 & .6925 & 7.8 \\
\bottomrule
\end{tabular}
\caption{First $40$ CLEVR reserved images. Destroying cross-patch timing collapses the
partition (\textsc{fg-ari}, \textsc{obj-iou}) while barely moving foreground detection, which
is rate-based. Substituting the phase gates for the spikes changes nothing at all.}
\label{tab:timing}
\end{table}

\paragraph{Gates in place of spikes.}
The third row replaces \texttt{component\_spikes} with \texttt{component\_gates} and changes
nothing else. Every mask is reproduced exactly, on all $40$ images and to every digit of all
three metrics. The reason is visible in the traces: with $s_{ia}=g_ie_{ia}$, the firing
indicator $e$ equals $1$ in $99.992\%$ of settled steps, the correlation between $s$ and $g$ is
$0.9997$, and the $0.008\%$ of steps where the neuron is silent never change a thresholded
decision. The timing that the shift ablation shows to be necessary is therefore the timing of
the oscillatory gate. The spiking layer carries it without loss, but the readout cannot tell
the two apart, and no claim should be made for the threshold nonlinearity on this evidence.

\subsection{Limitations}
\label{sec:limitations}

\paragraph{Object count.}
Our readout groups patches better than it counts objects. With the minimum-region-size rule it
predicts $5.15$ objects per image on CLEVR against a true $6.36$, matching the exact count on
$25.0\%$ of images; $2.64$ against $2.49$ on Multi-dSprites ($53.1\%$); and $5.28$ against exactly
$3$ on Tetrominoes ($9.4\%$). Without the rule the same model predicts $8.04$, $6.62$ and $21.07$,
matching on $17.5\%$, $28.1\%$ and \emph{none} of the $320$ Tetrominoes images, so most of the
count error was single-patch fragments rather than genuine over-segmentation. Letting $f$ be
chosen on each dataset's own validation split instead of transferred from CLEVR removes most of
the remainder --- $f=0.10$ gives $3.02$ predicted objects on Tetrominoes, $86.9\%$ exact, and
raises foreground IoU further to $.8066$ --- but that uses target-domain labels, so we report the
transferred value as the headline.

Two cautions about this metric. First, Tetrominoes contains exactly three objects in every image,
so a model that always answers three is exactly right on every image without grouping anything:
the fixed-$k$ baselines all score $100\%$ there, CtCAE included at \textsc{fg-ari} $0.0756$. The
count is only informative where the true count varies, and on Multi-dSprites, where it does, the
same baselines score $0\%$ exact against our $53.1\%$. Second, the residual error is a predicted
consequence of the objective: the model is trained without a reconstruction term, and a
reconstruction loss is what penalises splitting one object into two parts that must then be
reconstructed independently. Every baseline in Table~\ref{tab:main} optimises such a term; ours
does not, which is also what makes it cheap to train.

\paragraph{Seeds.}
Every number reported here is a single training seed. We measured the effect of reseeding one
baseline, which moved \textsc{fg-ari} by up to $0.134$ on CLEVR and $0.009$ on Multi-dSprites. All
margins in Table~\ref{tab:main} exceed that range, but this is evidence about baseline variance
and not an error bar on our own model.

\paragraph{Protocol.}
The patch protocol is not comparable with pixel-protocol numbers reported in the cited works, and
the held-out split is reserved from the training corpus rather than an independent benchmark test
set. Figures from other papers appear in this section only as each model's comparison with itself
(the baseline-fairness comparison), never as entries in Table~\ref{tab:main}.
